\documentclass[12pt,a4paper,oneside]{book} 

%%%%%%%%%%%%%%%%%%%%%%%%%%%%%%%%%%%%%%%%%%%%%%%%%%%%%%%%%%%%%%%%%%%%%%%%

\usepackage{ifthen}
\usepackage{graphicx}
\usepackage[dvips]{epsfig}
\usepackage{enumerate}
\usepackage{calc}
\usepackage{multicol} 
\usepackage{titlesec}
%\usepackage{showkeys}

%%%%%%%%%%%%%%%%%%%%%%%%%%%%%%%%%%%%%%%%%%%%%%%%%%%%%%%%%%%%%%%%%%%%%%%%

\usepackage{a4}
\usepackage{amsfonts}
\usepackage{amssymb}
\usepackage{epsfig}

%%%%%%%%%%%%%%%%%%%%%%%%%%%%%%%%%%%%%%%%%%%%%%%%%%%%%%%%%%%%%%%%%%%%%%%%

\usepackage{t1enc,times}
\usepackage{latexsym,amssymb}
\usepackage{amsmath}
\usepackage{amstext}
%\usepackage[T1]{fontenc}
%\usepackage{cmbright}
\usepackage{pifont}
\usepackage{marvosym}
%\usepackage{pslatex}
%\usepackage{stmaryrd}
%\usepackage{txfonts}

%%%%%%%%%%%%%%%%%%%%%%%%%%%%%%%%%%%%%%%%%%%%%%%%%%%%%%%%%%%%%%%%%%%%%%%%%%%%%%%

\usepackage{fancybox} 
\usepackage{fancyhdr}
\usepackage{fullpage}

%%%%%%%%%%%%%%%%%%%%%%%%%%%%%%%%%%%%%%%%%%%%%%%%%%%%%%%%%%%%%%%%%%%%%%%%%%%%%%%

%\usepackage{url}    
%\ExecuteOptions{dvips}
%\usepackage[pdftex,colorlinks=true]{hyperref}
%\hypersetup{backref,pdfpagemode=UseThumbs,pdfstartview=Fit,
%pdfpagelayout=SinglePage,pdfstartpage=1,colorlinks=true,menucolor=msc,
%anchorcolor=msc,pagecolor=msc,urlcolor=rfr,breaklinks=true,hyperfootnotes=true}

%%%%%%%%%%%%%%%%%%%%%%%%%%%%%%%%%%%%%%%%%%%%%%%%%%%%%%%%%%%%%%%%%%%%%%%%%%%%%%%

\pagestyle{fancy}
\fancyhf{} 
%\renewcommand{\headrulewidth}{1pt}
%\renewcommand{\footrulewidth}{1pt}
\renewcommand{\headwidth}{\textwidth}
\fancyhead[LE]{\leftmark}
\fancyhead[RO]{\small \rightmark}
\fancyfoot[C]{\thepage}

%%%%%%%%%%%%%%%%%%%%%%%%%%%%%%%%%%%%%%%%%%%%%%%%%%%%%%%%%%%%%%%%%%%%%%%%%%%%%%%

\tolerance 4000
\textwidth 17.00cm
\topmargin -0.30cm
\oddsidemargin -0.25cm
\evensidemargin -0.25cm
\textheight 23.00cm
%\headsep 12pt
\headheight 15pt
\footskip 60pt
%\parindent 12pt

%%%%%%%%%%%%%%%%%%%%%%%%%%%%%%%%%%%%%%%%%%%%%%%%%%%%%%%%%%%%%%%%%%%%%%%%%%%%%%%

\renewcommand{\rmdefault}{ptm}  % times
\renewcommand{\rmdefault}{phv}  % helvetica
\renewcommand{\rmdefault}{pbk}  % bookman
\renewcommand{\rmdefault}{ppl}  % palatino
\renewcommand{\sfdefault}{phv}  % helvetica as sans serif
\renewcommand{\ttdefault}{pcr}  % courier as fixed width
\renewcommand{\tabcolsep}{8pt}
\renewcommand{\arraystretch}{1.25}

%%%%%%%%%%%%%%%%%%%%%%%%%%%%%%%%%%%%%%%%%%%%%%%%%%%%%%%%%%%%%%%%%%%%%%%%%%%%%%%

\def\nn{\nonumber}
\def\f{{\frac}}
\def\pa{{\partial}}
\def\d{{\rm d}}
\def\l{\left}
\def\r{\right}
\def\Mpl{M_{_{\rm Pl}}}

%%%%%%%%%%%%%%%%%%%%%%%%%%%%%%%%%%%%%%%%%%%%%%%%%%%%%%%%%%%%%%%%%%%%%%%%%%%%%%%

\def\done{\marginpar {\scriptsize DONE}}
\def\check{\marginpar {\scriptsize CHECK}}

%%%%%%%%%%%%%%%%%%%%%%%%%%%%%%%%%%%%%%%%%%%%%%%%%%%%%%%%%%%%%%%%%%%%%%%%%%%%%%%

\begin{document}

%%%%%%%%%%%%%%%%%%%%%%%%%%%%%%%%%%%%%%%%%%%%%%%%%%%%%%%%%%%%%%%%%%%%%%%%%%%%%%%

\baselineskip 20pt

%%%%%%%%%%%%%%%%%%%%%%%%%%%%%%%%%%%%%%%%%%%%%%%%%%%%%%%%%%%%%%%%%%%%%%%%%%%%%%%

\pagenumbering{roman}

%%%%%%%%%%%%%%%%%%%%%%%%%%%%%%%%%%%%%%%%%%%%%%%%%%%%%%%%%%%%%%%%%%%%%%%%%%%%%%%

\thispagestyle{empty}
\topskip 15pt
\hrule\hrule\hrule\hrule\hrule
\vskip 20pt
\centerline{\Huge \bf Accelerating universe and} 
\vskip 15pt
\centerline{\Huge \bf the nature of dark energy} 
\vskip 20pt
\hrule\hrule\hrule\hrule\hrule
\vskip 30pt
\centerline{\Large A project report}
\vskip 8pt
\centerline{\Large submitted in partial fulfillment 
for the award of the degree of}
\vskip 8pt
\centerline{\Large Master of Science}
\vskip 8pt
\centerline{\Large in}
\vskip 8pt 
\centerline{\Large Physics}
\vskip 8pt
\centerline{\Large by}
\vskip 8pt
\centerline{\Large \bf Tirthankar Banerjee}
\vskip 8pt
\centerline{\Large under the guidance of}
\vskip 8pt
\centerline{\Large  Dr.~L.~Sriramkumar}
\vskip 30pt 
\begin{center}
\epsfig{file=iitm.ps, width=3.0cm, height=3.0cm}
\end{center}
\vskip 8pt 
\centerline{\Large \bf Department of Physics}
\vskip 8pt 
\centerline{\Large \bf Indian Institute of Technology Madras}
\vskip 8pt 
\centerline{\Large \bf Chennai~600036, India}
\vskip 8pt
\centerline{\Large \bf April~2012}
%%%%%%%%%%%%%%%%%%%%%%%%%%%%%%%%%%%%%%%%%%%%%%%%%%%%%%%%%%%%%%%%%%%%%%%%%%%%%%%

\newpage\topskip 40pt
\centerline{\Large CERTIFICATE}
\thispagestyle{empty}
\vskip 20pt\noindent 
This is to certify that the project titled {\bf On the generation of
gravitational waves during inflation} is a bona fide record of work done by 
{\bf Poruri Sai Rahul} towards the partial fulfillment of the 
requirements of the Master of Science degree ($?$) in Physics at the Indian 
Institute of Technology, Madras, Chennai 600036, India.
\vskip 120pt
\hspace{240pt}(L.~Sriramkumar, Project supervisor)

%%%%%%%%%%%%%%%%%%%%%%%%%%%%%%%%%%%%%%%%%%%%%%%%%%%%%%%%%%%%%%%%%%%%%%%%%%%%%%%

\newpage\topskip 40pt
\thispagestyle{empty}
\centerline{\Large ACKNOWLEDGEMENTS}
\vskip 20pt\noindent 
add stuff here

%%%%%%%%%%%%%%%%%%%%%%%%%%%%%%%%%%%%%%%%%%%%%%%%%%%%%%%%%%%%%%%%%%%%%%%%%%%%%%%

\newpage\topskip 40pt
\thispagestyle{empty}
\centerline{\Large ABSTRACT}
\vskip 20pt\noindent 
add stuff here

%%%%%%%%%%%%%%%%%%%%%%%%%%%%%%%%%%%%%%%%%%%%%%%%%%%%%%%%%%%%%%%%%%%%%%%%%%%%%%%

\newpage
\thispagestyle{empty}
\tableofcontents
\newpage

%%%%%%%%%%%%%%%%%%%%%%%%%%%%%%%%%%%%%%%%%%%%%%%%%%%%%%%%%%%%%%%%%%%%%%%%%%%%%%%

\pagenumbering{arabic}

%%%%%%%%%%%%%%%%%%%%%%%%%%%%%%%%%%%%%%%%%%%%%%%%%%%%%%%%%%%%%%%%%%%%%%%%%%%%%%%

\chapter{Introduction}

stuff~\cite{cite}.

\par 

more stuff

%%%%%%%%%%%%%%%%%%%%%%%%%%%%%%%%%%%%%%%%%%%%%%%%%%%%%%%%%%%%%%%%%%%%%%%%%%%%%%%

\section{Standard candles, and the role of supernovae}

Consider a source which emits light of a given wavelength, say, $\lambda_{\rm
e}$.
By the time the light reaches us, the universe would have expanded.
As we shall illustrate, a consequence of the expansion is the fact that the
wavelengths of light from other galaxies are stretched, and what we actually 
observe is light of a different wavelength, $\lambda_{\rm o}$. 
This phenomenon is known as redshift, which is characterized by the redshift 
parameter, $z= (\lambda_{\rm o}-\lambda_{\rm e})/\lambda_{\rm e}$. 
It so happens that the measurement of the brightness of far away galaxies as 
a function of the redshift proves to be as a very effective tool in
understanding 
the composition and, hence, the dynamics of the universe.

\par

One should appreciate the fact the redshift is an easily measurable quantity.
While the redshift associated with a single spectra line will, evidently, be 
impossible to determine, the presence of a wide variety of elements (and at
different levels of ionization) results in a pattern of a large number of 
lines in the spectra of galaxies. It is the complete pattern that facilities 
the determination of the redshift of galaxies.

\par

In contrast, the distance to an astronomical object is tricky to arrive at.
While methods such as trigonometric parallax are sturdy and straightforward,
their utility turns out to be limited when the distances involved are rather 
great and, hence, the angles involved very small.  
It is in this context that the concept of a standard candle comes in useful
(for an extended discussion on this topic, see, for instance,
Ref.~\cite{webb-1999}).
By a standard candle one refers to an astronomical object with a known
intrinsic 
luminosity or brightness.
Given the intrinsic luminosity, the observed brightness then allows us to 
estimate the distance to the object.
A simple example can illustrate the point.
If one is informed that a light bulb emits $100\; {\rm W}$ of energy, then
measuring the flux at our location allows us to arrive at the distance to 
the bulb from us.
Supernovae are exploding stars, which turn extraordinarily bright at their
peak intensity that they often outshine their host galaxy.
Certain types of supernovae---known as type Ia---exhibit such strong universal 
characteristics (which is related to the way they were formed) that their 
intrinsic, maximum brightness is very closely related to the pattern of the 
rise and fall of their intensities with time. 
The universal characteristic of these type Ia supernovae allow them to be used
as standard candles.
Moreover, their high peak brightness admit us to utilize them to determine
distances over cosmological scales.
Various observational missions such as the Supernova Cosmology 
Project (SCP)~\cite{scp} and the Super Nova Legacy Survey (SNLS)~\cite{snls} 
have assembled an impressive amount of data, which have led to interesting 
implications for our understanding of the universe.

%%%%%%%%%%%%%%%%%%%%%%%%%%%%%%%%%%%%%%%%%%%%%%%%%%%%%%%%%%%%%%%%%%%%%%%%%%%%%%%

\section{}


%%%%%%%%%%%%%%%%%%%%%%%%%%%%%%%%%%%%%%%%%%%%%%%%%%%%%%%%%%%%%%%%%%%%%%%%%%%%%%%

\section{Notations and conventions}


We shall set $\hbar=c=1$, but shall display $G$ explicitly, and define 
the Planck mass to be $\Mpl=\left(8\, \pi\, G\right)^{-1/2}$.

We shall work in $(3 + 1)$-spacetime dimensions, and adopt the metric 
signature of~$(+,-,-,-)$.
The Greek and Latin indices shall denote the spacetime and the spatial 
coordinates, respectively, 


%%%%%%%%%%%%%%%%%%%%%%%%%%%%%%%%%%%%%%%%%%%%%%%%%%%%%%%%%%%%%%%%%%%%%%%%%%%%%%%

\chapter{The Friedmann universe}



%%%%%%%%%%%%%%%%%%%%%%%%%%%%%%%%%%%%%%%%%%%%%%%%%%%%%%%%%%%%%%%%%%%%%%%%%%%%%%%

\section{}

\section{}

%%%%%%%%%%%%%%%%%%%%%%%%%%%%%%%%%%%%%%%%%%%%%%%%%%%%%%%%%%%%%%%%%%%%%%%%%%%%%%%

\chapter{Distant supernovae and the accelerating universe}

%%%%%%%%%%%%%%%%%%%%%%%%%%%%%%%%%%%%%%%%%%%%%%%%%%%%%%%%%%%%%%%%%%%%%%%%%%%%%%%

\chapter{Summary}

%%%%%%%%%%%%%%%%%%%%%%%%%%%%%%%%%%%%%%%%%%%%%%%%%%%%%%%%%%%%%%%%%%%%%%%%%%%%%%%

\begin{thebibliography}{99}
%%%%%%%%%%%%%%%%%%%%%%%%%%%%%%%%%%%%%%%%%%%%%%%%%%%%%%%%%%%%%%%%%%%%%%%%%%%%%%%
\bibitem{2df}
See, {\tt http://magnum.anu.edu.au/~TDFgg/}.
%%%%%%%%%%%%%%%%%%%%%%%%%%%%%%%%%%%%%%%%%%%%%%%%%%%%%%%%%%%%%%%%%%%%%%%%%%%%%%%
\bibitem{sdss}
See, {\tt http://www.sdss.org/}.
%%%%%%%%%%%%%%%%%%%%%%%%%%%%%%%%%%%%%%%%%%%%%%%%%%%%%%%%%%%%%%%%%%%%%%%%%%%%%%%
\bibitem{webb-1999}
S.~Webb, {\sl Measuring the Universe}\/ (Springer, Berlin, 1999).
%%%%%%%%%%%%%%%%%%%%%%%%%%%%%%%%%%%%%%%%%%%%%%%%%%%%%%%%%%%%%%%%%%%%%%%%%%%%%%%
\bibitem{scp}
See, {\tt http://supernova.lbl.gov/}.
%%%%%%%%%%%%%%%%%%%%%%%%%%%%%%%%%%%%%%%%%%%%%%%%%%%%%%%%%%%%%%%%%%%%%%%%%%%%%%%
\bibitem{snls}
See, {\tt http://snls.in2p3.fr/}.
%%%%%%%%%%%%%%%%%%%%%%%%%%%%%%%%%%%%%%%%%%%%%%%%%%%%%%%%%%%%%%%%%%%%%%%%%%%%%%%
\bibitem{durrer-2008}
R.~Durrer, {\sl The Cosmic Microwave Background}\/ (Cambridge University Press, 
Cambridge, England, 2008).
%%%%%%%%%%%%%%%%%%%%%%%%%%%%%%%%%%%%%%%%%%%%%%%%%%%%%%%%%%%%%%%%%%%%%%%%%%%%%%%
\bibitem{padmanabhan-1999}
T.~Padmanabhan, {\sl Theoretical Astrophysics, Volume III:~Galaxies and 
Cosmology},\/ (Cambridge University Press, Cambridge, England, 1999).
%%%%%%%%%%%%%%%%%%%%%%%%%%%%%%%%%%%%%%%%%%%%%%%%%%%%%%%%%%%%%%%%%%%%%%%%%%%%%%%
\end{thebibliography}
%%%%%%%%%%%%%%%%%%%%%%%%%%%%%%%%%%%%%%%%%%%%%%%%%%%%%%%%%%%%%%%%%%%%%%%%%%%%%%%
\end{document}